\documentclass[aspectratio=169,10pt,fontset=fandol]{ctexbeamer}
\usepackage{amsmath,amssymb,mathtools,bm,mathrsfs}
\usepackage{graphicx,booktabs,array,tabularx,multirow}
\usepackage{tikz,pgfplots,xparse,xstring,mdframed}
\usetikzlibrary{arrows.meta,calc,angles,positioning,shapes.geometric,graphs}
\pgfplotsset{compat=1.18}
\usetheme{fduclassroom}
\input{workflow_style.tex}
\input{workflow_fonts.tex}
\newenvironment{proofdetail}{\par\noindent\textbf{证明：}\quad}{\par}
\renewenvironment{solution}{\par\noindent\textbf{解：}\quad}{\par}
\title{中文 路径(样例).tex}
\begin{document}
\section{数列极限}
% FRAME-ID F001
\begin{frame}[t,fragile]{目录}
\normalsize
{\color{blue}\normalsize 1. 数列极限\par}\vspace{0.3em}
{\color{gray!40}\normalsize 2. 极限运算\par}\vspace{0.3em}
\end{frame}
% FRAME-ID F002
\begin{frame}[t,fragile]{数列极限}
\normalsize
\begin{block}{定义}
{\normalsize
% SOURCE-BEGIN U0001 0 77
若对任意 $\varepsilon>0$，存在 $N$，当 $n>N$ 时 $|a_n-a|<\varepsilon$，则称 $a_n$ 收敛于 $a$。
% SOURCE-END U0001
}
\end{block}
\end{frame}
% FRAME-ID F003
\begin{frame}[t,fragile]{数列极限}
\normalsize
\normalsize
% ADDITION-BEGIN b3c8ec39e5164535abfefc0b3f71a7b2 S01
补充例题：求常数数列的极限。
% ADDITION-END b3c8ec39e5164535abfefc0b3f71a7b2
\end{frame}
\section{极限运算}
% FRAME-ID F004
\begin{frame}[t,fragile]{目录}
\normalsize
{\color{black}\normalsize 1. 数列极限\par}\vspace{0.3em}
{\color{blue}\normalsize 2. 极限运算\par}\vspace{0.3em}
\end{frame}
% FRAME-ID F005
\begin{frame}[t,fragile]{极限运算}
\normalsize
{\normalsize
% SOURCE-BEGIN U0004 0 49

若 $a_n\to a$ 且 $b_n\to b$，则 $a_n+b_n\to a+b$。

% SOURCE-END U0004
}
\end{frame}
\end{document}
